% Source: https://artofproblemsolving.com/community/c3h2359190_which_is_it_closest_to
% Stable identity: aops-topic-2359190-q05
Take decimal logarithms in Stirling's formula with $n=2020$:
\[\log_{10}2020!=2020\left(\log_{10}2020-\log_{10}e\right)+\dfrac12\log_{10}(2\pi\cdot2020)+\theta\log_{10}e.\]
With $\log_{10}2020\approx3.3053514$ and $\log_{10}e\approx0.4342945$, the first term is $2020\cdot2.8710569\approx5799.5349$; the second is $\dfrac12\log_{10}12692.0\approx2.0518$; the last is less than $0.00002$. So
\[\log_{10}2020!\approx5801.587,\qquad 2020!\approx3.86\times10^{5801}.\]

\textbf{(a)} On a logarithmic scale $2020!$ is within a factor of $3$ of $10^{5802}$ ($2020!$ has $5802$ digits), while every other choice differs from it by a factor of more than $10^{1800}$. Answer: $\boxed{\textbf{(C)}}$.

\textbf{(b)} The absolute differences are:
\begin{itemize}
\item to (C): $10^{5802}-2020!\approx(10-3.86)\times10^{5801}=6.14\times10^{5801}$;
\item to (B): $2020!-10^{4000}<2020!\approx3.86\times10^{5801}$;
\item to (A): about $3.86\times10^{5801}$ as well, but larger than for (B);
\item to (D) and (E): more than $10^{7999}$.
\end{itemize}
So the closest in absolute terms is $\boxed{\textbf{(B) }10^{4000}}$. The reason is that $2020!$ is less than half of $10^{5802}$, so it is nearer to $0$, and hence to the tiny number $10^{4000}$, than to $10^{5802}$.
